% typeup: the look of typed-up study notes. Loaded with
% pandoc --include-in-header (the typeup CLI does this) together with
% marks.lua. The mark boxes are what marks.lua emits for
% blockquotes that start with **[Added]**, **[Your mark]**, **[Unclear]**
% or another **[...]** source label.
\usepackage[most]{tcolorbox}
\usepackage{xcolor}

% Tables: alternate rows get a light grey band and a little padding, so
% cells that wrap onto several lines stay visibly apart from the next row.
% \rowcolors needs xcolor's table option: typeup passes it as a class
% option, because the pandoc template loads xcolor before this header.
\definecolor{rowstripe}{RGB}{242,242,242}
\usepackage{etoolbox}
\AtBeginEnvironment{longtable}{%
  \rowcolors{2}{}{rowstripe}%
  \renewcommand{\arraystretch}{1.25}}

\definecolor{addedframe}{RGB}{30,90,170}
\definecolor{addedback}{RGB}{235,242,252}
\definecolor{yourframe}{RGB}{190,120,0}
\definecolor{yourback}{RGB}{253,245,225}
\definecolor{unclearframe}{RGB}{190,30,45}
\definecolor{unclearback}{RGB}{253,232,234}
\definecolor{sourceframe}{RGB}{110,110,110}
\definecolor{sourceback}{RGB}{243,243,243}
\definecolor{metacol}{RGB}{110,110,110}

% In the "How to read this note" legend (notelegend, below) each mark box
% says after its name what it means; elsewhere the name stands alone
\newtoggle{inlegend}
\newcommand{\legendrole}[1]{\iftoggle{inlegend}{\enspace{\mdseries\color{metacol}#1}}{}}
% What each mark means, as the legend says it. A note's legend-added,
% legend-your, legend-unclear and legend-source metadata (or the config
% file's [legend] table) replace these (marks.lua, docs/config.md)
\newcommand{\addedrole}{added while typing up: check it against the source}
\newcommand{\yourrole}{a note the writer made in the margin}
\newcommand{\unclearrole}{an uncertain reading: check it against the original}
\newcommand{\sourcerole}{from the source named in the label}

% Common box style: thin left rule, coloured title on the same line
\tcbset{notemark/.style={enhanced, breakable, boxrule=0pt, leftrule=2.5pt,
  arc=0pt, outer arc=0pt, left=6pt, right=6pt, top=3pt, bottom=3pt,
  before skip=6pt, after skip=6pt, fontupper=\small,
  attach title to upper={\enspace}, fonttitle=\small\bfseries}}

\newtcolorbox{markadded}{notemark, colframe=addedframe, colback=addedback,
  coltitle=addedframe, title={Added\legendrole{\addedrole}}}
\newtcolorbox{markyour}{notemark, colframe=yourframe, colback=yourback,
  coltitle=yourframe, title={Your mark\legendrole{\yourrole}}}
\newtcolorbox{markunclear}{notemark, colframe=unclearframe,
  colback=unclearback, coltitle=unclearframe, title={Unclear\legendrole{\unclearrole}}}
% Any other **[...]** label: the label text is the title
\newtcolorbox{marksource}[1]{notemark, colframe=sourceframe,
  colback=sourceback, coltitle=sourceframe, title={#1\legendrole{\sourcerole}}}

% "How to read this note": the ::: legend div at the top of a note, with
% one real, shortened example of each mark the note uses (marks.lua)
\newtcolorbox{notelegend}{enhanced, boxrule=0.4pt, arc=1.5pt,
  colframe=metacol!40, colback=white, left=7pt, right=7pt, top=4pt,
  bottom=4pt, before skip=8pt, after skip=10pt, fontupper=\small,
  title={How to read this note}, fonttitle=\small\bfseries,
  coltitle=metacol, colbacktitle=white, toptitle=3pt, bottomtitle=0pt,
  before upper={\toggletrue{inlegend}\tcbset{notemark/.append style={
    before skip=3pt, after skip=5pt, breakable=false}}}}

% Inline [x_i?]{.unclear}: a reading to check against the page
\newcommand{\unclearinline}[1]{%
  \tcbox[on line, boxrule=0pt, arc=1pt, boxsep=0pt, left=2pt, right=2pt,
    top=1pt, bottom=1pt, colback=unclearback, coltext=unclearframe]{#1}}

% Content boxes: what KIND of text it is (the mark boxes above say where
% it came from). marks.lua turns a fenced div into one of these:
%   ::: {.definition title="Dual problem"}  ...  :::
% Definition, Theorem, Proposition, Lemma, Corollary, Example, Remark and
% Exercise share one counter, so "Definition 3" is followed by "Example 4"
% and any number is easy to find. Proof has no number and ends with a box.
% The counter restarts with each note in a combined PDF.
\definecolor{nbdefframe}{RGB}{0,115,105}
\definecolor{nbthmframe}{RGB}{105,60,150}
\definecolor{nbexframe}{RGB}{55,125,55}
\definecolor{nbremframe}{RGB}{75,90,110}
\definecolor{nbexerframe}{RGB}{160,65,115}
\definecolor{nbproofframe}{RGB}{150,150,150}

\tcbset{notebox/.style={enhanced, breakable, boxrule=0.5pt, arc=1.5pt,
  outer arc=1.5pt, left=7pt, right=7pt, top=4pt, bottom=4pt,
  toptitle=2pt, bottomtitle=2pt, before skip=8pt, after skip=8pt,
  fonttitle=\bfseries\small}}

% \nbbox{env}{Label}{colour}: a numbered box; its argument is the optional
% title in brackets after the number ({} for none)
\newcommand{\nbtitle}[2]{#1~\thetcbcounter\ifblank{#2}{}{\enspace(#2)}}
\newtcolorbox[auto counter]{nbdefinition}[1]{notebox,
  colframe=nbdefframe, colback=nbdefframe!5, colbacktitle=nbdefframe!15,
  coltitle=nbdefframe!70!black, title={\nbtitle{Definition}{#1}}}
\newcommand{\nbbox}[3]{%
  \newtcolorbox[use counter from=nbdefinition]{#1}[1]{notebox,
    colframe=#3, colback=#3!5, colbacktitle=#3!15, coltitle=#3!70!black,
    title={\nbtitle{#2}{##1}}}}
\nbbox{nbtheorem}{Theorem}{nbthmframe}
\nbbox{nbproposition}{Proposition}{nbthmframe}
\nbbox{nblemma}{Lemma}{nbthmframe}
\nbbox{nbcorollary}{Corollary}{nbthmframe}
\nbbox{nbexample}{Example}{nbexframe}
\nbbox{nbremark}{Remark}{nbremframe}
\nbbox{nbexercise}{Exercise}{nbexerframe}
% Proof: no number, a thin grey rule on the left, a box at the end
\newtcolorbox{nbproof}[1]{notemark, colframe=nbproofframe, colback=white,
  coltitle=black, fontupper=\normalsize, fonttitle=\itshape\small,
  title={\ifblank{#1}{Proof.}{Proof (#1).}},
  after upper={\par\nobreak\hfill$\square$}}

% Figures (plots made with the typeup.notes matplotlib style, docs/figures.md): stay where
% they are written instead of floating to another page, centred, with a
% small grey caption under them
\usepackage{float}
\makeatletter\def\fps@figure{H}\makeatother
\usepackage{caption}
\DeclareCaptionFont{meta}{\color{metacol}}
\captionsetup[figure]{font={small,meta}, labelfont={bf,meta},
  skip=4pt, belowskip=4pt}

% Keep-with-next (marks.lua puts \needspace before headings and
% lead-in lines, so they never end a page alone)
\usepackage{needspace}

% Let TeX stretch spaces a little more before a line sticks out
\setlength{\emergencystretch}{2em}

% ---- Page furniture (ideas from the Eisvogel pandoc template, BSD-3) ----
% Which note this is: marks.lua sets these from the metadata of a single
% note (course, label, date), and
% \notebreak sets them per note in a combined PDF.
\newcommand{\notecourse}{}
\newcommand{\notelabel}{}
\newcommand{\notedate}{}
\definecolor{headcol}{RGB}{25,55,95}

% Running header (course . lecture on the left, date on the right) and a
% "page n / N" footer on every page; the title page of a combined PDF
% has the footer only
\usepackage{fancyhdr}
\usepackage{lastpage}
\setlength{\headheight}{14pt}
\pagestyle{fancy}
\fancyhf{}
\renewcommand{\headrulewidth}{0.4pt}
\renewcommand{\headrule}{{\color{metacol!50}\hrule height\headrulewidth}}
\fancyhead[L]{\footnotesize\color{metacol}\notecourse
  \ifdefempty{\notelabel}{}{\enspace\textperiodcentered\enspace\notelabel}}
\fancyhead[R]{\footnotesize\color{metacol}\notedate}
\fancyfoot[C]{\footnotesize\color{metacol}\thepage\ / \pageref*{LastPage}}
\fancypagestyle{plain}{\fancyhf{}\renewcommand{\headrulewidth}{0pt}%
  \fancyfoot[C]{\footnotesize\color{metacol}\thepage\ / \pageref*{LastPage}}}

% Headings: the note's title (#) in navy with a rule under it, then the
% course . lecture . date line that marks.lua puts after the first #
\usepackage{titlesec}
\titleformat{\section}{\Large\bfseries\color{headcol}}{}{0pt}{}
  [{\vspace{2pt}\color{headcol!60}\titlerule[0.8pt]}]
\titleformat{\subsection}{\large\bfseries\color{headcol}}{}{0pt}{}
\titleformat{\subsubsection}{\normalsize\bfseries\color{headcol!80!black}}{}{0pt}{}
\titlespacing*{\section}{0pt}{4pt}{10pt}
\titlespacing*{\subsection}{0pt}{14pt}{6pt}
\titlespacing*{\subsubsection}{0pt}{10pt}{4pt}
\newcommand{\notemetaline}{%
  \ifdefempty{\notecourse}{}{%
    {\noindent\small\color{metacol}\notecourse
     \ifdefempty{\notelabel}{}{\enspace\textperiodcentered\enspace\notelabel}%
     \ifdefempty{\notedate}{}{\enspace\textperiodcentered\enspace\notedate}\par}%
    \vspace{6pt}}}

% Code blocks: a light grey frame, for highlighted code (pandoc's Shaded)
% and plain ``` blocks alike (marks.lua wraps those in codebox)
\definecolor{codeframe}{RGB}{200,200,200}
\definecolor{codeback}{RGB}{248,248,248}
\newtcolorbox{codebox}{enhanced, breakable, boxrule=0.4pt, arc=1.5pt,
  colframe=codeframe, colback=codeback, left=6pt, right=6pt, top=3pt,
  bottom=3pt, before skip=6pt, after skip=8pt, fontupper=\small}
\makeatletter
\AtBeginDocument{\@ifundefined{Shaded}{}{%
  \renewenvironment{Shaded}{\begin{codebox}}{\end{codebox}}}}

% Each note in a combined PDF starts on a new page (typeup combine),
% restarts the box numbers and sets its own header:
%   \notebreak{course}{label}{date}
\newcommand{\notebreak}[3]{\clearpage\setcounter{tcb@cnt@nbdefinition}{0}%
  \renewcommand{\notecourse}{#1}\renewcommand{\notelabel}{#2}%
  \renewcommand{\notedate}{#3}}
\makeatother
